\documentclass[11pt]{amsart}

\usepackage{amsmath}
\usepackage{amssymb}
\usepackage{booktabs}
\usepackage{array}
\usepackage[margin=1.15in]{geometry}
\usepackage{url}
\usepackage[hidelinks]{hyperref}

\allowdisplaybreaks

\newtheorem{theorem}{Theorem}[section]
\newtheorem{proposition}[theorem]{Proposition}
\newtheorem{lemma}[theorem]{Lemma}
\newtheorem{corollary}[theorem]{Corollary}
\theoremstyle{definition}
\newtheorem{definition}[theorem]{Definition}
\newtheorem{example}[theorem]{Example}
\newtheorem{problem}[theorem]{Problem}
\theoremstyle{remark}
\newtheorem{remark}[theorem]{Remark}

\numberwithin{equation}{section}
\emergencystretch=2em

\newcommand{\repo}[1]{\texttt{\detokenize{#1}}}
% The public repository of this paper.
\newcommand{\repourl}{https://github.com/no-way-labs/smale6-generic-finiteness}
\newcommand{\QQ}{\mathbb{Q}}
\newcommand{\Qbar}{\overline{\mathbb{Q}}}
\newcommand{\CC}{\mathbb{C}}
\newcommand{\RR}{\mathbb{R}}
\newcommand{\ZZ}{\mathbb{Z}}
\newcommand{\cN}{\mathcal{N}}
\newcommand{\cX}{\mathcal{X}}
\newcommand{\cM}{\mathcal{M}}
\newcommand{\Gm}{\mathbb{G}}
\newcommand{\Ps}{\Psi^{*}}
\newcommand{\Es}{E^{*}}
\newcommand{\Po}{\Psi^{\circ}}
\newcommand{\Eo}{E^{\circ}}
\newcommand{\one}[1]{\mathbf{1}_{#1}}
\newcommand{\nd}[2]{[#1]_{#2}}
\newcommand{\step}[1]{\smallskip\noindent\textit{#1}\ }
\DeclareMathOperator{\Spec}{Spec}

\begin{document}

\title[Generic finiteness of planar central configurations]{Generic finiteness of planar central configurations for every number of bodies}

\author{Keston Aquino-Michaels}
\address{No Way Labs}

\begin{abstract}
We prove that for every $n \ge 2$ there is a proper Zariski-closed subset of $\CC^n$, defined over
$\QQ$, such that for positive masses outside it the planar $n$-body problem has finitely many central
configurations up to similarity. The subset is not described, and we attempt no count of the
configurations sharper than the bound of B\'ezout's inequality. The statement is known for $n \le 5$ and is the subject of two preprints of 2026 for
$n = 6$; we know no earlier proof for $n \ge 7$. We work with the system of Albouy and Kaloshin over a
field of Hahn series, at masses whose valuations $u_0 < u_1 < \dots < u_{n-1}$ satisfy
$u_j - u_{j-1} > 2n^2 (u_{j-1} - u_0)$. At such masses we show that the least valuation among the terms
of the potential and of the moment of inertia is one of finitely many numbers fixed by the $u_j$, and
hence that the valuation of the square of every mutual distance of a solution, which measures how
closely two bodies approach, is bounded. A continuum of solutions modulo rotation would contain
solutions in which this valuation is as large as one likes. The proof contains no
enumeration of diagrams and no computation. Formal counterparts of the statements about valuations are
proved in Lean~4 with Mathlib: that the valuations of a solution satisfy the cut conditions
(Lemma~\ref{lem:A}), Lemma~\ref{lem:B}, Lemma~\ref{lem:C}, Theorem~\ref{thm:D4}(a), (b), (c) and the
inequality of~(d), and Corollary~\ref{cor:bound}. The reduction to these statements, which uses algebraic geometry over the
field of Hahn series, is not formalised.
\end{abstract}

\maketitle

\section{Introduction}\label{sec:intro}

Let $n \ge 2$. Pairwise distinct points $z_0, \dots, z_{n-1} \in \CC$ with masses $m_0, \dots, m_{n-1}$
form a \emph{central configuration} if there are $\lambda \in \RR$ and $c \in \CC$ such that
\begin{equation}\label{eq:cc}
  \lambda\,(z_t - c) \;=\; \sum_{l \ne t} m_l\, \frac{z_t - z_l}{|z_t - z_l|^{3}}
  \qquad (t = 0, \dots, n-1).
\end{equation}
Central configurations are the shapes of the homographic solutions of the Newtonian $n$-body problem,
and in particular of its relative equilibria~\cite{wintner1941}. A \emph{similarity} is a map
$z \mapsto \alpha z + \beta$ or $z \mapsto \alpha \bar z + \beta$ with $\alpha \ne 0$; it maps a central
configuration to a central configuration. The number of classes is finite with the reflections if and
only if it is finite without them. The question whether, for every choice of positive masses, the
number of central configurations up to similarity is finite goes back to Wintner~\cite{wintner1941} and
is the sixth problem on Smale's list~\cite{smale1998}.

For three bodies there are five classes up to similarities without reflection, for all positive masses:
the three collinear configurations of Euler and the two equilateral triangles of Lagrange.
Table~\ref{tab:known} lists what is proved or claimed for $n \ge 4$. For four bodies Hampton and Moeckel proved finiteness for all positive masses.
For five bodies Albouy and Kaloshin proved it outside an explicit subvariety of codimension~2 of the mass
space, and three preprints of October 2026 give proofs for all positive masses. For six bodies two preprints
of 2026 give proofs of finiteness for positive masses outside a proper algebraic subset: that of Tomar, whose
proof is formalised in Lean, and Part~I of a work of Chen and Tong, which builds on the classification of
diagrams by Chang and Chen~\cite{changchen2024jsc,changchen2025siads}. Tomar's proof contains its own
enumeration of the diagrams of Albouy and Kaloshin for six bodies, and both preprints exclude the diagrams
with computer algebra. For seven bodies Moczurad
and Zgliczy\'nski list the central configurations with equal masses and count those with masses close
to equal, with computer-assisted proofs.

\begin{table}[ht]
\caption{Finiteness of planar central configurations for $n \ge 4$ bodies: results and claims, with the
hypothesis on the masses and the state of the source. The table rests on a search of arXiv and Zenodo
made on 17 November 2026.}\label{tab:known}
\small
\begin{tabular}{@{}l>{\raggedright}p{0.30\textwidth}>{\raggedright}p{0.27\textwidth}>{\raggedright\arraybackslash}p{0.25\textwidth}@{}}
\toprule
$n$ & masses & source & state \\
\midrule
4 & all positive & Hampton and Moeckel~\cite{hamptonmoeckel2006} & published \\
5 & positive, outside an explicit subvariety of codimension 2 & Albouy and Kaloshin~\cite{albouykaloshin2012} & published \\
5 & outside a proper algebraic subset & Jensen and Leykin~\cite{jensenleykin2025} & published; a computation of a tropical prevariety \\
5 & all positive & Tomar~\cite{tomar5}; Chen and Tong~\cite{chentong5}; Liu and Liu~\cite{liuliu5} & three preprints \\
6 & positive, outside a proper algebraic subset & Tomar~\cite{tomar6} & preprint; the proof is formalised in Lean \\
6 & positive, outside a proper algebraic subset & Chen and Tong~\cite{chentong6} & preprint; rests on the classification of Chang and Chen and on computational certificates; the exceptional masses are left to a Part~II \\
6 & cross configurations with $m_5 = m_6$, positive masses outside a proper algebraic subset of these & Dias and Pan~\cite{diaspan2020} & published \\
5, 6, 7 & equal masses & Moczurad and Zgliczy\'nski~\cite{moczuradzgliczynski2019} & published; computer-assisted; all configurations are listed \\
4 to 7 & a neighbourhood of the equal masses in the usual topology & Moczurad and Zgliczy\'nski~\cite{moczuradzgliczynski2020} & published; computer-assisted; the configurations are counted \\
all & Dziobek configurations (of dimension $n - 2$, planar only for $n = 4$), outside a proper algebraic subset & Moeckel~\cite{moeckel2001} & published \\
all & some open sets of positive masses in the usual topology & Xia~\cite{xia1991} & published; the configurations are counted \\
all & all positive & Tibboel~\cite{tibboel2013} & withdrawn by the author \\
all & not stated & Okolo~\cite{okolo2026} & only the title is public \\
\bottomrule
\end{tabular}
\end{table}

No entry of Table~\ref{tab:known} gives finiteness outside a proper algebraic subset of the mass space
for any $n \ge 7$. The results of Xia and of Moczurad and Zgliczy\'nski hold at the equal masses or on
open sets of masses in the usual topology, and say nothing outside them. Of the two claims for every $n$, one was withdrawn and the other is a record whose
abstract and file are not public. The following theorem gives finiteness outside a proper algebraic
subset for every $n$.

\begin{theorem}\label{thm:main}
For every $n \ge 2$ there is a proper Zariski-closed subset $\Sigma_n$ of $\CC^n$, defined over $\QQ$,
with the following property. For every $(m_0, \dots, m_{n-1}) \in \RR_{>0}^n$ that does not lie in
$\Sigma_n$, the points of the plane that form a central configuration with the masses
$m_0, \dots, m_{n-1}$ fall into finitely many classes up to similarity.
\end{theorem}

The proof is about a system of polynomial equations with complex unknowns. For masses $m_0, \dots, m_{n-1}$ in
a field $k$, the system $\cN_n$ has the unknowns $X_t$, $Y_t$ for $0 \le t < n$ and $P_{tl} = P_{lt}$ for
$0 \le t < l < n$. With $X_{tl} = X_t - X_l$ and $Y_{tl} = Y_t - Y_l$ its equations are
\begin{align}
  P_{tl}^2\, X_{tl}\, Y_{tl} &= 1 && (t \ne l), \label{eq:N1}\\
  X_t &= \sum_{l \ne t} m_l\, P_{tl}^3\, X_{tl} && (0 \le t < n), \label{eq:N2}\\
  Y_t &= \sum_{l \ne t} m_l\, P_{tl}^3\, Y_{tl} && (0 \le t < n). \label{eq:N3}
\end{align}
These are the equations of Albouy and Kaloshin~\cite{albouykaloshin2012}: a central configuration with
$c = 0$ and $\lambda = 1$ gives the solution $X_t = z_t$, $Y_t = \bar z_t$, $P_{tl} = |z_t - z_l|^{-1}$.
An element $s \in k^{*}$ maps a solution $(X, Y, P)$ to the solution $(sX, s^{-1}Y, P)$. We refer to
this action of $k^{*}$ as the \emph{rotation}, because for $k = \CC$ and $|s| = 1$ it is a rotation of
the plane.

\begin{theorem}\label{thm:Nn}
For every $n \ge 2$ there is a proper Zariski-closed subset $\Sigma_n$ of $\CC^n$, defined over $\QQ$,
such that for every $m \in \CC^n$ outside $\Sigma_n$ the system $\cN_n$ with the masses $m$ has finitely
many solutions in $\CC$ modulo the rotation.
\end{theorem}

Let $K = \Qbar((T^{\RR}))$ be the field of Hahn series with coefficients in $\Qbar$ and real exponents,
with its valuation $v$ (Section~\ref{sec:val}). For $\tau \ge 0$ let
\begin{equation}\label{eq:chamber}
  W_n(\tau) = \bigl\{\, u \in \RR^n : u_1 > u_0 \text{ and } u_j - u_{j-1} > \tau\,(u_{j-1} - u_0)
  \text{ for } 2 \le j \le n-1 \,\bigr\}.
\end{equation}

\begin{theorem}\label{thm:K}
Let $n \ge 2$ and let $m \in (K^{*})^n$ be such that $(v(m_0), \dots, v(m_{n-1})) \in W_n(2n^2)$. Then the
system $\cN_n$ with the masses $m$ has finitely many solutions in $K$ modulo the rotation.
\end{theorem}

Theorem~\ref{thm:main} follows from Theorem~\ref{thm:Nn} and Lemmas~\ref{lem:norm}
and~\ref{lem:cc-sol}. Theorem~\ref{thm:Nn} follows from Theorem~\ref{thm:K} and
Lemma~\ref{lem:D0}. Theorem~\ref{thm:K} follows from Lemma~\ref{lem:A}, Corollary~\ref{cor:bound},
Proposition~\ref{prop:D2} and Lemma~\ref{lem:orbits}. Corollary~\ref{cor:bound} is a consequence of
Theorem~\ref{thm:D4}, which uses Lemmas~\ref{lem:B} and~\ref{lem:C}.

\begin{remark}\label{rem:scope}
The set $\Sigma_n$ is not described, and we attempt no count of the solutions beyond the following one.
When the number of solutions of $\cN_n$ modulo the rotation is finite, it is at most
$4^{2n + \binom{n}{2}}$: the solutions with $X_{01} = 1$ contain one point of each orbit, and they are
the common zeros of $2n + \binom{n}{2}$ polynomials of degree at most $4$ and one of degree $1$, to
which B\'ezout's inequality~\cite{heintz1983} applies. In
Theorem~\ref{thm:Nn} the set $\Sigma_n$ cannot avoid the positive masses: for $n = 5$ it contains
$(4, 4, 4, 4, 1)$, where $\cN_5$ has a continuum of solutions that are not central configurations
(Remark~\ref{rem:roberts}). Whether $\Sigma_n$ can be chosen disjoint from $\RR_{>0}^n$ in
Theorem~\ref{thm:main} is Smale's problem, and this paper does not address it.
\end{remark}

\subsection*{Method} Albouy and Kaloshin~\cite{albouykaloshin2012} classify by diagrams the ways in which
a sequence of solutions can become singular, and exclude the diagrams one by one. The six-body
preprints~\cite{tomar6,chentong6} follow this approach for the positive central configurations, with
computer algebra, and Jensen and Leykin~\cite{jensenleykin2025} compute the recession cones of a tropical
prevariety at chosen valuations of the masses in a field of Puiseux series. The number of diagrams and the size of the computation grow with $n$. The
present proof has no classification, and three of its steps are taken from the literature. The
system $\cN_n$ is system~(4) of~\cite[\S 2]{albouykaloshin2012} in the coordinates $X_t = x_t + i y_t$,
$Y_t = x_t - i y_t$, with $P_{tl}$ in place of their $\delta_{tl}$ and without their normalisation
$y_{12} = 0$.
Proposition~\ref{prop:D2} is the valuation form of an argument of theirs: a polynomial on a closed
algebraic set takes finitely many values or all values but finitely
many~\cite[Lemma~1]{albouykaloshin2012}, and on a continuum of solutions modulo rotation some mutual
distance is of the second kind (the proofs of Theorems~3 and~6
of~\cite{albouykaloshin2012}). The proof of Lemma~\ref{lem:D0} embeds the field of the rational functions of the
masses in $K$ in such a way that the masses receive prescribed valuations, and Theorem~\ref{thm:K} is a
theorem at those valuations, where~\cite{jensenleykin2025} has a computation. The identity obtained when the
equations \eqref{eq:N2} are added over a subset of the bodies underlies the proof of Rule~1c
of~\cite[p.~546]{albouykaloshin2012}, and Tomar~\cite[\S 2.3, \S 4.1 and \S 9.1]{tomar6} states it for
both colours and uses it with the origin as an additional point of both trees of clusters.

The step that we did not find in the literature is Theorem~\ref{thm:D4}. Let the masses have valuations
in $W_n(2n^2)$, so that every body is lighter than the previous one by a large power of $T$. At a solution
over $K$, consider the terms $m_t m_l P_{tl}$ of the potential and the terms $m_t X_t Y_t$ of the moment
of inertia, and among them the set $\Es$ of the terms of least valuation. Each term of $\Es$ has a cluster in the
tree of the $X$ coordinates and a cluster in the tree of the $Y$ coordinates. The identities of
Lemma~\ref{lem:A} show that each such cluster belongs to at least two terms of $\Es$, and therefore the
graph that joins the two clusters of each term has a cycle. Around a cycle the heights of the clusters
cancel, and what remains is a relation $N \Ps + q \cdot u = 0$ between the least valuation $-\Ps$ and the
valuations $u$ of the masses, with integers $N$ and $q_j$ bounded in terms of $n$. Lemma~\ref{lem:C}
shows that a cycle cannot consist of terms of the potential alone, and a parity argument, which uses the
exponent $3$ of the Newtonian force, then gives $N \ne 0$. Hence $\Ps$ is one of finitely many numbers
fixed by $u$, and every mutual distance of every solution has a valuation bounded in terms of $u$.

\subsection*{Formal proof and computations} Formal counterparts of the following statements are proved
in Lean~4 with Mathlib: the last sentence of Lemma~\ref{lem:A}, Lemma~\ref{lem:B}, Lemma~\ref{lem:C} and
Corollary~\ref{cor:C}, Theorem~\ref{thm:D4}(a), (b), (c) and the inequality of~(d), and
Corollary~\ref{cor:bound}. They are proved against three files of statements, each of which a separate
session compared with the text without access to the formal proofs. Table~\ref{tab:lean} in
Section~\ref{sec:lean} says what each formal statement covers. Lemmas~\ref{lem:norm}
to~\ref{lem:orbits}, Proposition~\ref{prop:D2} and Lemma~\ref{lem:D0} are not formalised, and neither are
the deductions of Theorems~\ref{thm:main}, \ref{thm:Nn} and~\ref{thm:K} from them. Floating-point searches
for five and six bodies agree with Theorem~\ref{thm:D4}, and two exact searches agree with
Lemma~\ref{lem:C} for at most six bodies (Section~\ref{sec:comp}); they are evidence, and no proof in this
paper uses them. The proofs, the formal proofs and the programs were produced with a language-model agent,
and Section~\ref{sec:llm} records how.

\subsection*{Organisation} Section~\ref{sec:system} relates central configurations to the system
$\cN_n$. Section~\ref{sec:val} introduces the valuations, the two trees and the cut identities.
Sections~\ref{sec:chamber} to~\ref{sec:D4} contain the statements about valuations: Lemma~\ref{lem:B},
Lemma~\ref{lem:C} and Theorem~\ref{thm:D4}. Section~\ref{sec:K} proves Theorem~\ref{thm:K} and
Section~\ref{sec:generic} proves Theorems~\ref{thm:Nn} and~\ref{thm:main}. Section~\ref{sec:hyp} says
where each hypothesis is used, Sections~\ref{sec:lean} and~\ref{sec:comp} describe the formal proof and
the computations, Section~\ref{sec:problems} states open problems, and Section~\ref{sec:llm} records the
use of language models.

\section{Central configurations and the system \texorpdfstring{$\cN_n$}{N\_n}}\label{sec:system}

Throughout, $n \ge 2$, the \emph{bodies} are the indices $0, \dots, n-1$, and $t$, $l$, $a$, $b$, $c$
denote bodies. A central configuration is \emph{normalised} if \eqref{eq:cc} holds with $c = 0$ and
$\lambda = 1$.

\begin{lemma}\label{lem:norm}
If the masses are positive, every central configuration is similar to a normalised one. If moreover two
normalised central configurations are similar, then one is the image of the other under a map
$z \mapsto e^{i\theta} z$ or $z \mapsto e^{i\theta} \bar z$ with $\theta \in \RR$.
\end{lemma}

\begin{proof}
Multiply \eqref{eq:cc} by $m_t\, \overline{(z_t - c)}$ and add over $t$. On the right, the two terms of a
pair $\{t, l\}$ add up to $m_t m_l / |z_t - z_l|$. This gives $\lambda I = U$ with
$I = \sum_t m_t |z_t - c|^2$ and $U = \sum_{t < l} m_t m_l / |z_t - z_l|$. Both $I$ and $U$ are positive,
and therefore $\lambda > 0$. The translation by $-c$ gives $c = 0$, and the map
$z \mapsto \lambda^{1/3} z$ then gives $\lambda = 1$, because a scaling by $\rho > 0$ divides $\lambda$
by $\rho^{3}$. For the second statement, add \eqref{eq:cc} multiplied by $m_t$ over $t$: the right side
vanishes, and hence $c$ is the centre of mass. A similarity $z \mapsto \alpha z + \beta$ or
$z \mapsto \alpha \bar z + \beta$ between two normalised configurations preserves the centre of mass
$0$, which gives $\beta = 0$, and it divides $\lambda$ by $|\alpha|^{3}$, which gives $|\alpha| = 1$.
\end{proof}

\begin{lemma}\label{lem:cc-sol}
Let $m \in \RR_{>0}^n$.
\begin{enumerate}
\item[(a)] If $(z_t)$ is a normalised central configuration, then $X_t = z_t$, $Y_t = \bar z_t$,
$P_{tl} = |z_t - z_l|^{-1}$ is a solution of $\cN_n$ in $\CC$.
\item[(b)] If the solutions of two normalised central configurations $(z_t)$ and $(z'_t)$ lie in one orbit
of the rotation, then $z'_t = s z_t$ for all $t$ with one $s$ of modulus $1$.
\item[(c)] If $\cN_n$ has finitely many solutions in $\CC$ modulo the rotation, then the central
configurations with the masses $m$ fall into finitely many classes up to similarity.
\end{enumerate}
\end{lemma}

\begin{proof}
For (a), $X_{tl} Y_{tl} = |z_t - z_l|^2$ gives \eqref{eq:N1}, and \eqref{eq:N2} is \eqref{eq:cc} with
$c = 0$ and $\lambda = 1$; \eqref{eq:N3} is its complex conjugate. For (b), let $z'_t = s z_t$ and
$\bar z'_t = s^{-1} \bar z_t$ for all $t$. The conjugate of the second equation is
$z'_t = \bar s^{-1} z_t$, and some $z_t$ is not $0$, which gives $s \bar s = 1$. For (c), the map of (a)
sends the normalised central configurations modulo rotations of the plane injectively to the orbits, by
(b). By Lemma~\ref{lem:norm} every similarity class contains a normalised central configuration.
\end{proof}

\begin{lemma}\label{lem:facts}
Let $(X, Y, P)$ be a solution of $\cN_n$ over a field $k$.
\begin{enumerate}
\item[(a)] For all $t \ne l$ the elements $X_{tl}$, $Y_{tl}$ and $P_{tl}$ are not $0$. At most one $X_t$
is $0$, and at most one $Y_t$ is $0$.
\item[(b)] If $s \in k^{*}$ fixes the solution, then $s = 1$.
\end{enumerate}
\end{lemma}

\begin{proof}
Part (a) follows from \eqref{eq:N1}, and for (b) the equation $s X_{01} = X_{01}$ with $X_{01} \ne 0$
gives $s = 1$.
\end{proof}

\begin{lemma}\label{lem:orbits}
Let $k$ be an algebraically closed field and $m \in k^n$. The set of the solutions of $\cN_n$ in $k$ with
the masses $m$ is a Zariski-closed subset of an affine space over $k$. It consists of finitely many orbits
of the rotation if and only if none of its irreducible components has dimension at least $2$.
\end{lemma}

\begin{proof}
The group $k^{*}$ is connected, and therefore it maps every irreducible component to itself; by
Lemma~\ref{lem:facts}(b) every orbit has dimension $1$. An orbit that lies in a component of dimension
$1$ is a constructible subset of that component of the same dimension; it therefore contains a nonempty
open subset of the component. Two nonempty open subsets of an irreducible space meet, and hence a
component of dimension $1$ is a single orbit. If no component has dimension at least $2$, every
component is one orbit, and there are finitely many components. Conversely, a finite union of orbits has
dimension at most $1$.
\end{proof}

\section{Valuations, trees and the cut identities}\label{sec:val}

\subsection{The field of Hahn series} Let $K = \Qbar((T^{\RR}))$ be the field of the formal sums
$\sum_{\gamma} a_\gamma T^{\gamma}$ with $a_\gamma \in \Qbar$, real exponents $\gamma$ and a well-ordered
support. The valuation $v(x)$ of $x \ne 0$ is the least exponent of $x$, and $v(0) = +\infty$. The field
$K$ is algebraically closed~\cite[Theorem~1 and Corollary~4]{poonen1993}, and we use three properties
of its valuation. First, $v(xy) = v(x) + v(y)$ and $v(x + y) \ge \min(v(x), v(y))$. Second, if finitely many elements
of $K$, not all $0$, have the sum $0$, then the least of their valuations is the valuation of at least
two of them. Third, every real number is the valuation of infinitely many elements of $K$.

\subsection{Ultrametrics and nodes}

\begin{definition}\label{def:ultra}
An \emph{ultrametric} on a finite set $\Lambda$ is a symmetric function
$(x, y) \mapsto d_{xy}$ on the pairs of distinct elements of $\Lambda$, with values in
$\RR \cup \{+\infty\}$, such that $d_{xz} \ge \min(d_{xy}, d_{yz})$ for all distinct $x, y, z$; we put
$d_{xx} = +\infty$. For $x \ne y$ with $d_{xy}$ real, the \emph{node} of the pair $\{x, y\}$ is the set
$\nd{x, y}{d} = \{ p \in \Lambda : d_{xp} \ge d_{xy} \}$, and $d_{xy}$ is its \emph{height}. The
\emph{children} of a node $A$ of height $h$ are the classes of the relation ``$p = p'$ or $d_{pp'} > h$''
on $A$.
\end{definition}

The condition on $d$ says that in every triangle the least of the three values is taken at least twice.
With the convention $d_{xx} = +\infty$ the inequality $d_{xz} \ge \min(d_{xy}, d_{yz})$ holds for all
$x, y, z \in \Lambda$, distinct or not, and the proofs below use it in this form.

\begin{lemma}\label{lem:tree}
Let $d$ be an ultrametric on $\Lambda$, let $x \ne y$ have a real value $h = d_{xy}$, and let
$A = \nd{x, y}{d}$.
\begin{enumerate}
\item[(a)] The node $A$ contains $x$ and $y$, it equals $\nd{y, x}{d}$, and $d_{pp'} \ge h$ for all
$p, p' \in A$. If $\nd{x', y'}{d} = A$ for a pair with a real value $d_{x'y'}$, then $d_{x'y'} = h$.
\item[(b)] The relation of Definition~\ref{def:ultra} is an equivalence relation on $A$, the elements $x$
and $y$ lie in different children, and $d_{pp'} = h$ whenever $p$ and $p'$ lie in different children.
\item[(c)] If $p \in A$ and $p' \notin A$, then $d_{pp'} < h$.
\item[(d)] If $p, p' \in A$ and $d_{pp'} = h$, then $\nd{p, p'}{d} = A$.
\end{enumerate}
\end{lemma}

\begin{proof}
(a) If $d_{xp} \ge h$ then $d_{yp} \ge \min(d_{yx}, d_{xp}) \ge h$, and the converse holds by symmetry;
this gives $\nd{x, y}{d} = \nd{y, x}{d}$. For $p, p' \in A$ we have
$d_{pp'} \ge \min(d_{px}, d_{xp'}) \ge h$. If $\nd{x', y'}{d} = A$, then $x', y' \in A$ gives
$d_{x'y'} \ge h$, and $x, y \in \nd{x', y'}{d}$ gives $h \ge d_{x'y'}$. (b) The relation is transitive
because $d_{pp''} \ge \min(d_{pp'}, d_{p'p''})$. The elements $x$ and $y$ are not related because
$d_{xy} = h$. If $p$ and $p'$ lie in different children, then $d_{pp'} \le h$ by the definition of the
relation and $d_{pp'} \ge h$ by (a). (c) Here $d_{xp'} < h \le d_{xp}$. If $d_{pp'}$ were at least $h$,
then $d_{xp'} \ge \min(d_{xp}, d_{pp'}) \ge h$. (d) Let $q \in \Lambda$. Because
$d_{xp} \ge h$, the inequality $d_{pq} \ge \min(d_{px}, d_{xq})$ shows that $d_{xq} \ge h$ implies
$d_{pq} \ge h$, and the inequality $d_{xq} \ge \min(d_{xp}, d_{pq})$ shows the converse. Hence
$\nd{p, p'}{d} = \{ q : d_{pq} \ge h \} = A$.
\end{proof}

\subsection{The weights of a solution} Let $m \in (K^{*})^n$ and put $u_t = v(m_t)$, the \emph{mass
weight} of the body $t$; a larger mass weight is a lighter body. For a solution of $\cN_n$ over $K$ put
\[
  \xi_t = v(X_t), \qquad \eta_t = v(Y_t), \qquad \zeta_{tl} = v(X_{tl}), \qquad \omega_{tl} = v(Y_{tl}).
\]
By Lemma~\ref{lem:facts}(a) the numbers $\zeta_{tl}$ and $\omega_{tl}$ are real, and $\xi_t$, $\eta_t$
are real or $+\infty$; the equation \eqref{eq:N1} gives $v(P_{tl}) = -(\zeta_{tl} + \omega_{tl})/2$.

We add to the $n$ bodies one more \emph{vertex} $O$, the origin, and put $\zeta_{tO} = \xi_t$ and
$\omega_{tO} = \eta_t$. Then $\zeta_{xy}$ is the valuation of the difference of two of the $n + 1$
elements $X_0, \dots, X_{n-1}, 0$ of $K$, and therefore $\zeta$ is an ultrametric on the $n + 1$
vertices; the same holds for $\omega$. The letters $x$, $y$, $p$ denote vertices, and an \emph{edge} is
a pair of vertices: a \emph{body pair} $\{t, l\}$ or an \emph{$O$-edge} $\{t, O\}$. We speak of the
\emph{colour $z$} for $X$, $\zeta$, $\xi$ and of the \emph{colour $w$} for $Y$, $\omega$, $\eta$. For an
edge $e = \{x, y\}$ we write $\zeta_e = \zeta_{xy}$ and $\omega_e = \omega_{xy}$, and, when these are
real, $\nd{e}{\zeta}$ and $\nd{e}{\omega}$ for its two nodes.

\begin{definition}\label{def:psi}
Let $u \in \RR^n$ and let $\zeta$, $\omega$ be ultrametrics on the $n + 1$ vertices with real values on
the body pairs. Put $\xi_t = \zeta_{tO}$ and $\eta_t = \omega_{tO}$, and
\[
  \Psi(t, l) = \tfrac12 (\zeta_{tl} + \omega_{tl}) - u_t - u_l, \qquad
  \Psi(t, O) = -\xi_t - \eta_t - u_t ,
\]
where $\Psi(t, O) = -\infty$ if $\xi_t$ or $\eta_t$ is $+\infty$. Let $\Ps$ be the largest value of
$\Psi$ on the edges, a real number because $n \ge 2$, and let $\Es$ be the set of the edges $e$ with
$\Psi(e) = \Ps$. For a nonempty set $S$ of bodies, the \emph{terms of the cut of $S$ in the colour $z$}
are the body pairs $\{t, l\}$ with $t \in S$ and $l \notin S$, with the values
$\Phi^z(t, l) = \Psi(t, l) + \omega_{tl}$, and the $O$-edges $\{t, O\}$ with $t \in S$ and $\xi_t$ real,
with the values $\Phi^z(t, O) = -u_t - \xi_t$. The \emph{cut condition} for $S$ in the colour $z$ says
that two different terms of the cut have the same value $\Phi^z$ and no term has a larger one. The terms
and the cut condition in the colour $w$ are defined in the same way with $\zeta$ and $\omega$ exchanged
and $\eta$ in place of $\xi$: $\Phi^w(t, l) = \Psi(t, l) + \zeta_{tl}$ and $\Phi^w(t, O) = -u_t - \eta_t$.
The triple $(u, \zeta, \omega)$ is a \emph{weight datum} if the cut condition holds for every nonempty
set of bodies and both colours.
\end{definition}

For the weights of a solution, $-\Psi(t, l)$ is the valuation of the term $m_t m_l P_{tl}$ of the
potential $\sum_{t < l} m_t m_l P_{tl}$, and $-\Psi(t, O)$ is the valuation of the term $m_t X_t Y_t$ of
the moment of inertia $\sum_t m_t X_t Y_t$. If $\xi_t$ and $\eta_t$ are real, then
$\Phi^z(t, O) = \Psi(t, O) + \eta_t$ and $\Phi^w(t, O) = \Psi(t, O) + \xi_t$, which is the rule
$\Phi^z(e) = \Psi(e) + \omega_e$ of the body pairs.

\begin{lemma}[cut identities]\label{lem:A}
Let $m \in (K^{*})^n$ and let $(X, Y, P)$ be a solution of $\cN_n$ over $K$. For every nonempty set $S$
of bodies,
\begin{equation}\label{eq:cut}
  \sum_{t \in S} m_t X_t \;-\; \sum_{t \in S,\ l \notin S} m_t m_l\, P_{tl}^3\, X_{tl} \;=\; 0 ,
\end{equation}
and the same identity holds with $Y$ in place of $X$. The term of a body pair $\{t, l\}$ in
\eqref{eq:cut} has the valuation $-\Phi^z(t, l)$, and the term $m_t X_t$ has the valuation
$-\Phi^z(t, O)$ if $X_t \ne 0$. The weights $(u, \zeta, \omega)$ of the solution form a weight datum.
\end{lemma}

\begin{proof}
Multiply \eqref{eq:N2} for the body $t$ by $m_t$ and add over $t \in S$. A pair $\{t, l\}$ inside $S$
contributes $m_t m_l P_{tl}^3 (X_{tl} + X_{lt}) = 0$, which gives \eqref{eq:cut}. The valuation of
$m_t m_l P_{tl}^3 X_{tl}$ is
$u_t + u_l - \tfrac32 (\zeta_{tl} + \omega_{tl}) + \zeta_{tl} = -(\Psi(t, l) + \omega_{tl})$, and the
valuation of $m_t X_t$ is $u_t + \xi_t$. The terms of \eqref{eq:cut} that are not $0$ are therefore the
terms of the cut of $S$ in the colour $z$, with the valuations $-\Phi^z$. At least one of them is not
$0$: a set $S$ that does not contain all bodies has the term of a body pair, and the set of all bodies
has a term $m_t X_t \ne 0$ by Lemma~\ref{lem:facts}(a) and $n \ge 2$. By the second property of $v$, the
least valuation among the terms that are not $0$ is taken by two of them. This is the cut condition in
the colour $z$, and the colour $w$ is the same argument with \eqref{eq:N3}.
\end{proof}

\begin{example}[two bodies]\label{ex:n2}
For $n = 2$ the system gives $P_{01}^3 (m_0 + m_1) = 1$. If $u_0 < u_1$, then
$\zeta_{01} + \omega_{01} = 2u_0/3$, $\xi_1 = \zeta_{01}$ and $\xi_0 = \zeta_{01} + u_1 - u_0$, and the
same holds for $\eta$ and $\omega$. Hence $\Psi(0, 1) = \Psi(1, O) = -(2u_0 + 3u_1)/3$, and
$\Psi(0, O)$ is smaller by $u_1 - u_0$. The set $\Es$ consists of $\{0, 1\}$ and $\{1, O\}$, and both
edges have the node $\{0, 1, O\}$ in both colours.
\end{example}

The identity \eqref{eq:cut} underlies the proof of Rule~1c of~\cite[p.~546]{albouykaloshin2012}, a
statement about the centre of mass of an isolated component of a diagram.
Tomar~\cite[\S 2.3, \S 4.1 and \S 9.1]{tomar6} states the identity for both colours and uses it with the
vertex $O$ in both trees. Table~\ref{tab:notation} collects the notation of Sections~\ref{sec:system}
to~\ref{sec:D4}.

\begin{table}[ht]
\caption{Notation.}\label{tab:notation}
\small
\begin{tabular}{@{}>{\raggedright}p{0.30\textwidth}>{\raggedright\arraybackslash}p{0.64\textwidth}@{}}
\toprule
symbol & meaning \\
\midrule
$t, l, a, b, c$; $x, y, p$; $e, f, g$ & bodies; vertices (bodies and $O$); edges \\
$\cN_n$; $X_t, Y_t, P_{tl}$; $X_{tl}, Y_{tl}$ & the system \eqref{eq:N1} to \eqref{eq:N3}; its unknowns; $X_t - X_l$, $Y_t - Y_l$ \\
$K$, $v$, $T$ & the field of Hahn series, its valuation, its variable \\
$u_t$ & the mass weight $v(m_t)$ \\
$\xi_t, \eta_t$; $\zeta_{xy}, \omega_{xy}$ & $v(X_t)$, $v(Y_t)$; the two ultrametrics on the vertices \\
$\nd{e}{\zeta}$, $\nd{e}{\omega}$ & the nodes of the edge $e$ (Definition~\ref{def:ultra}) \\
$\Psi$, $\Ps$, $\Es$ & Definition~\ref{def:psi} \\
$\Phi^z$, $\Phi^w$ & the values of the terms of a cut (Definition~\ref{def:psi}) \\
$W_n(\tau)$ & the chamber \eqref{eq:chamber} \\
$\Po$, $\Eo$ & the largest $\Psi$ on the body pairs and the pairs that attain it (Section~\ref{sec:C}) \\
$f_i, g_i$; $\varepsilon$; $N$, $q$ & the edges of a cycle; their signs; Definition~\ref{def:cycle} \\
$\one{t}$; $\one{e}$ & the unit vector of the body $t$ in $\ZZ^n$; the sum of $\one{t}$ over the bodies $t$ of the edge $e$ \\
\bottomrule
\end{tabular}
\end{table}

\section{The chamber}\label{sec:chamber}

The set $W_n(\tau)$ of \eqref{eq:chamber} is an open cone. It is not empty: the vector with $u_0 = 0$
and $u_j = (\tau + 2)^{j-1}$ for $j \ge 1$ lies in it. For $\tau \ge 0$ and $u \in W_n(\tau)$,
induction on $j$ gives
\begin{equation}\label{eq:mono}
  0 < u_1 - u_0 < u_2 - u_0 < \dots < u_{n-1} - u_0
  \quad\text{and}\quad u_j - u_0 > (1 + \tau)(u_{j-1} - u_0) \ \text{ for } 2 \le j \le n-1 .
\end{equation}
In particular the mass weights increase with the index: the body $0$ is the heaviest.

\begin{lemma}\label{lem:B}
Let $n \ge 2$, $\tau \ge 0$ and $u \in W_n(\tau)$. Let $q \in \ZZ^n$ be a nonzero vector with
$\sum_j q_j = 0$, and let $p$ be the largest index with $q_p \ne 0$. If
\[
  \tau \;\ge\; \frac{1}{|q_p|} \sum_{1 \le j < p} |q_j| \;-\; 1 ,
\]
then $q \cdot u \ne 0$, and $q \cdot u$ has the sign of $q_p$.
\end{lemma}

\begin{proof}
Because $\sum_j q_j = 0$ and $q \ne 0$, we have $p \ge 1$ and
$q \cdot u = \sum_{1 \le j \le p} q_j (u_j - u_0)$. If $p = 1$ this is $q_1 (u_1 - u_0)$. Let $p \ge 2$.
By \eqref{eq:mono} and the hypothesis on $\tau$,
\[
  \Bigl| \sum_{1 \le j < p} q_j (u_j - u_0) \Bigr|
  \;\le\; \Bigl( \sum_{1 \le j < p} |q_j| \Bigr) (u_{p-1} - u_0)
  \;\le\; (1 + \tau)\, |q_p|\, (u_{p-1} - u_0) \;<\; |q_p|\, (u_p - u_0) ,
\]
and therefore the term $q_p (u_p - u_0)$ decides the sign of $q \cdot u$.
\end{proof}

\begin{remark}\label{rem:tau}
The hypothesis $\tau \ge 0$ is needed: the vector $u = (0, 2, 3/2, 1)$ satisfies the inequalities of
\eqref{eq:chamber} with $n = 4$ and $\tau = -1/2$, the vector $q = (-1, -1, 0, 2)$ satisfies the inequality of the lemma, and $q \cdot u = 0$.
\end{remark}

\section{Cycles of body pairs}\label{sec:C}

In this section the vertex $O$ does not occur. Let $\zeta$ and $\omega$ be ultrametrics on the $n$
bodies with real values, and let $u \in \RR^n$. With
$\Psi(a, b) = \tfrac12(\zeta_{ab} + \omega_{ab}) - u_a - u_b$ as in Definition~\ref{def:psi}, let $\Po$
be the largest value of $\Psi$ on the pairs of bodies and $\Eo$ the set of the pairs that attain it.

\begin{definition}\label{def:alt}
Let $E$ be a set of edges on which $\zeta$ and $\omega$ are real. A \emph{closed alternating sequence}
in $E$ is a sequence $f_0, g_0, f_1, g_1, \dots, f_{r-1}, g_{r-1}$ of $2r$ distinct edges of $E$, with
$r \ge 1$, such that $\omega_{f_i} = \omega_{g_i}$ and $\zeta_{g_i} = \zeta_{f_{i+1}}$ for
$0 \le i < r$, where $f_r = f_0$. It is a \emph{simple cycle} of $E$ if moreover
$\nd{f_i}{\omega} = \nd{g_i}{\omega}$ and $\nd{g_i}{\zeta} = \nd{f_{i+1}}{\zeta}$ for $0 \le i < r$,
the $r$ nodes $\nd{f_i}{\omega}$ are distinct, and the $r$ nodes $\nd{f_i}{\zeta}$ are distinct.
\end{definition}

A simple cycle of $E$ is a simple cycle, in the usual sense, of the following bipartite multigraph
$\Gamma(E)$. Its vertex set is the disjoint union of the set of the nodes $\nd{e}{\zeta}$, $e \in E$,
and the set of the nodes $\nd{e}{\omega}$, $e \in E$; a set of vertices that is a node in both colours
gives two vertices of $\Gamma(E)$. It has one edge for each $e \in E$, which joins $\nd{e}{\zeta}$ to
$\nd{e}{\omega}$. For $r = 1$ a simple cycle consists of two distinct edges of $E$ with the same node in
each colour. Two edges with the same node have the same height by Lemma~\ref{lem:tree}(a), and hence
the conditions on the nodes imply the conditions on the heights.

\begin{lemma}\label{lem:C}
Let $n \ge 2$, $\tau \ge n^2$ and $u \in W_n(\tau)$, and let $\zeta$ and $\omega$ be ultrametrics on
the $n$ bodies with real values. Let $\Eo$ be the set of the pairs of bodies at which
$\Psi(a, b) = \tfrac12(\zeta_{ab} + \omega_{ab}) - u_a - u_b$ takes its largest value. Then there is no
closed alternating sequence in $\Eo$.
\end{lemma}

\begin{corollary}\label{cor:C}
Under the hypotheses of Lemma~\ref{lem:C} the graph $\Gamma(\Eo)$ has no simple cycle.
\end{corollary}

\begin{proof}[Proof of Lemma~\ref{lem:C}]
Suppose that $f_0, g_0, \dots, f_{r-1}, g_{r-1}$ is a closed alternating sequence in $\Eo$. Step~1
shows that every body lies in as many pairs $f_i$ as pairs $g_i$. Step~2 expresses the heights along
the sequence by integer functions, one for each body. Steps~3 and~4 find, at the lightest body of the
sequence, three bodies whose triangle condition contradicts Lemma~\ref{lem:B}.

Put $\alpha_i = \zeta_{f_i}$ and $\beta_i = \omega_{f_i}$ for $0 \le i < r$, and $\alpha_r = \alpha_0$.
Then $\beta_i = \omega_{g_i}$ and $\alpha_{i+1} = \zeta_{g_i}$. All pairs of the sequence have the value
$\Po$, which gives, for $0 \le i < r$,
\begin{equation}\label{eq:rel}
  \alpha_i + \beta_i = 2\Po + 2\, \one{f_i} \cdot u
  \qquad\text{and}\qquad
  \alpha_{i+1} + \beta_i = 2\Po + 2\, \one{g_i} \cdot u .
\end{equation}

\step{Step 1: the vector $q = \sum_{i < r} (\one{g_i} - \one{f_i})$ is $0$.} The difference of the two
equations of \eqref{eq:rel} is $\alpha_{i+1} - \alpha_i = 2 (\one{g_i} - \one{f_i}) \cdot u$, and the
sum over $i$ is $0 = 2\, q \cdot u$. The coefficients of $q$ are integers with the sum $0$. A body lies
in at most $n - 1$ pairs and the $2r$ pairs are distinct, and therefore $|q_j| \le n - 1$ for every $j$.
If $q \ne 0$, let $p$ be its largest index with $q_p \ne 0$. The sum of Lemma~\ref{lem:B} has at most
$n - 2$ terms, it is at most $(n - 2)(n - 1)$, and $\tau \ge n^2$ exceeds this number; Lemma~\ref{lem:B}
then gives $q \cdot u \ne 0$, a contradiction. Hence $q = 0$, which says that every body lies in as
many pairs $g_i$ as pairs $f_i$.

\step{Step 2: levels.} For a body $k$ and $0 \le i \le r$ put
\[
  a_k(i) = \#\{ j < i : k \in g_j \} - \#\{ j < i : k \in f_j \} ,
\]
and for $0 \le i < r$ put $b_k(i) = a_k(i) - 1$ if $k \in f_i$ and $b_k(i) = a_k(i)$ otherwise. Then
$a_k(i + 1) = b_k(i) + 1$ if $k \in g_i$ and $a_k(i + 1) = b_k(i)$ otherwise, and
$a_k(r) = q_k = 0 = a_k(0)$ by Step~1. We claim that for $0 \le i \le r$ and $0 \le j < r$
\begin{equation}\label{eq:height}
  \alpha_i + \beta_j - 2\Po \;=\; 2 \sum_{k} \bigl( a_k(i) - b_k(j) \bigr)\, u_k .
\end{equation}
Indeed, the differences $\alpha_{i+1} - \alpha_i = 2 \sum_k (a_k(i+1) - a_k(i))\, u_k$ add up to
$\alpha_i = \alpha_0 + 2 \sum_k a_k(i)\, u_k$. The first equation of \eqref{eq:rel} then gives
$\beta_j = 2\Po - \alpha_0 - 2 \sum_k b_k(j)\, u_k$, and the sum of the two is \eqref{eq:height}. The
integers $a_k(i) - b_k(j)$ have three properties that Step~4 uses. Their sum over $k$ is $2$, because every pair
has two bodies, which gives $\sum_k a_k(i) = 0$ and $\sum_k b_k(j) = -2$. They are $0$ for a body $k$
that lies in no pair of the sequence. Their absolute value is at most $n$, because $a_k(i) - a_k(j)$ is
a difference of two numbers of distinct pairs through $k$, each between $0$ and $n - 1$.

\step{Step 3: two pairs through the lightest body.} Let $L$ be the largest index of a body that lies in
a pair of the sequence. Because $q_L = 0$, the body $L$ lies in some $f_i$. Consider the cyclic sequence
of the $2r$ integers $a_L(0), b_L(0), a_L(1), b_L(1), \dots, a_L(r-1), b_L(r-1)$. From $a_L(i)$ to
$b_L(i)$ it falls by $1$ if $L \in f_i$ and does not change otherwise; from $b_L(i)$ to $a_L(i+1)$ it
rises by $1$ if $L \in g_i$ and does not change otherwise. The sequence is not constant, and we let $\mu$ be its least value. Some entry equal to $\mu$ follows an entry larger than $\mu$, and this is a step from $a_L(i)$
to $b_L(i)$ with $L \in f_i$. Some entry equal to $\mu$ precedes an entry larger than $\mu$, and this is
a step from $b_L(j)$ to $a_L(j+1)$ with $L \in g_j$. Hence there are $i$ and $j$ with
\begin{equation}\label{eq:levels}
  L \in f_i, \quad L \in g_j, \quad b_L(i) = b_L(j) = \mu, \quad a_L(i) = a_L(j + 1) = \mu + 1 .
\end{equation}

\step{Step 4: the triangle.} Let $f_i = \{x, L\}$ and $g_j = \{x', L\}$ with $i$, $j$ as in
\eqref{eq:levels}. The pairs $f_i$ and $g_j$ are distinct, and therefore $x \ne x'$. Both $x$ and $x'$
have smaller indices than $L$; hence $L \ge 2$, and $u_x < u_L$ and $u_{x'} < u_L$ by \eqref{eq:mono}.
The ultrametric conditions for the three bodies $x$, $L$, $x'$ are
\[
  \zeta_{xx'} \ge \min(\zeta_{xL}, \zeta_{x'L}) = \min(\alpha_i, \alpha_{j+1})
  \qquad\text{and}\qquad
  \omega_{xx'} \ge \min(\omega_{xL}, \omega_{x'L}) = \min(\beta_i, \beta_j) ,
\]
and $\Psi(x, x') \le \Po$ says $\zeta_{xx'} + \omega_{xx'} \le 2\Po + 2u_x + 2u_{x'}$. Therefore
\begin{equation}\label{eq:four}
  \alpha_c + \beta_d \;\le\; 2\Po + 2u_x + 2u_{x'}
  \qquad\text{for some } c \in \{i, j + 1\} \text{ and some } d \in \{i, j\} .
\end{equation}
Each of the four cases of $(c, d)$ leads to a contradiction, as follows. If $(c, d) = (i, i)$, the first equation of \eqref{eq:rel} turns \eqref{eq:four}
into $u_L \le u_{x'}$, which is false. If $(c, d) = (j + 1, j)$, the second equation of \eqref{eq:rel}
turns it into $u_L \le u_x$, which is false. In the two other cases \eqref{eq:height} turns
\eqref{eq:four} into $\nu \cdot u \ge 0$ for the vector
\[
  \nu \;=\; \one{x} + \one{x'} - \sum_k \bigl( a_k(c) - b_k(d) \bigr)\, \one{k} \;\in\; \ZZ^n ,
  \qquad (c, d) = (i, j) \text{ or } (j + 1, i) .
\]
By Step~2 the coefficients of $\nu$ have the sum $0$, they are $0$ at the indices larger than $L$, and
$|\nu_k| \le n + 1$. By \eqref{eq:levels} the coefficient at $L$ is $\nu_L = -1$ in both cases. Hence
$L$ is the largest index with a nonzero coefficient, and the sum of Lemma~\ref{lem:B} for $\nu$ is at
most $(L - 1)(n + 1) \le (n - 2)(n + 1) < n^2 \le \tau$. Lemma~\ref{lem:B} therefore gives $\nu \cdot u < 0$, against $\nu \cdot u \ge 0$, and with this all
four cases are excluded.
\end{proof}

\begin{example}\label{ex:0011}
The hypothesis on $u$ is needed, as the following data for $n = 4$ and $u = (0, 0, 1, 1)$ show. Let
$\zeta_{ab} = 2$ if
$a, b \in \{1, 2, 3\}$ and $\zeta_{ab} = 0$ otherwise, and let $\omega_{ab} = 2$ if $a, b \in \{0, 2, 3\}$
and $\omega_{ab} = 0$ otherwise. Then $\Psi = 0$ on all six pairs, and
$\{0, 1\}, \{1, 2\}, \{2, 3\}, \{0, 3\}$ is a simple cycle of $\Eo$.
\end{example}

\section{The cycle relation}\label{sec:D4}

Let $(u, \zeta, \omega)$ be a weight datum, with $\Psi$, $\Ps$ and $\Es$ as in Definition~\ref{def:psi}.
An edge of $\Es$ has a real value $\Psi$, and therefore $\zeta$ and $\omega$ are real on $\Es$.

\begin{definition}\label{def:cycle}
For a closed alternating sequence $f_0, g_0, \dots, f_{r-1}, g_{r-1}$ in $\Es$ put
$\varepsilon(f_i) = -1$ and $\varepsilon(g_i) = +1$, and
\[
  N = 3 \sum_{e \text{ a body pair}} \varepsilon(e) \in \ZZ , \qquad
  q = \sum_{e \text{ a body pair}} 2\, \varepsilon(e)\, \one{e}
      \;-\; \sum_{e = \{t, O\}} \varepsilon(e)\, \one{t} \in \ZZ^n ,
\]
where the sums run over the edges $e$ of the sequence.
\end{definition}

\begin{theorem}\label{thm:D4}
Let $n \ge 2$, $\tau \ge 2n^2$ and $u \in W_n(\tau)$, and let $(u, \zeta, \omega)$ be a weight datum.
\begin{enumerate}
\item[(a)] For every edge $e \in \Es$ there is an edge $e' \in \Es$ with $e' \ne e$ and
$\nd{e'}{\omega} = \nd{e}{\omega}$, and an edge $e'' \in \Es$ with $e'' \ne e$ and
$\nd{e''}{\zeta} = \nd{e}{\zeta}$. The graph $\Gamma(\Es)$ has a simple cycle.
\item[(b)] Every closed alternating sequence in $\Es$ contains an $O$-edge.
\item[(c)] For every closed alternating sequence in $\Es$, the numbers of Definition~\ref{def:cycle}
satisfy $N \ne 0$ and $N \Ps + q \cdot u = 0$.
\item[(d)] The number $\Ps$ belongs to a finite set that depends on $n$ and $u$ only, and
$\Ps \le \tfrac13 (2n - 1) \sum_{j} |u_j|$.
\end{enumerate}
\end{theorem}

\begin{proof}
\step{Step 1: proof of (a).} Let $e = \{x, y\} \in \Es$, let $A = \nd{e}{\omega}$ and $h = \omega_e$. By
Lemma~\ref{lem:tree}(b) the vertices $x$ and $y$ lie in different children of $A$, and at most one of
these two children contains $O$. The pair $e$ is not ordered, and we name its vertices in such a way that the child $S$ of $A$ that
contains $x$ does not contain $O$. Then $S$ is a nonempty set of bodies and $y \notin S$. Let $g = \{p, p'\}$ be a term of the cut of $S$ in the colour $z$, with $p \in S$ and
$p' \notin S$. We show that
\begin{equation}\label{eq:phi}
  \Phi^z(g) = \Psi(g) + \omega_g \le \Ps + h, \quad\text{with equality only if } g \in \Es
  \text{ and } \nd{g}{\omega} = A .
\end{equation}
The equation $\Phi^z(g) = \Psi(g) + \omega_g$ holds for a body pair by definition. Let $g = \{p, O\}$,
so that $\xi_p$ is real. If $\eta_p = \omega_{pO}$ were $+\infty$, it would be larger than $h$; the
vertex $O$ would then lie in $A$, by Lemma~\ref{lem:tree}(c), and in the child $S$ of $p$. Hence
$\eta_p$ is real, and $\Phi^z(g) = \Psi(g) + \eta_p$ by the remark after Definition~\ref{def:psi}. For
the inequality, $\Psi(g) \le \Ps$, and $\omega_g \le h$: if $p' \in A$ then $p$ and $p'$ lie in different
children and $\omega_g = h$ by Lemma~\ref{lem:tree}(b), and if $p' \notin A$ then $\omega_g < h$ by
Lemma~\ref{lem:tree}(c). If equality holds, then $\Psi(g) = \Ps$ and $\omega_g = h$ with $p, p' \in A$,
and $\nd{g}{\omega} = A$ by Lemma~\ref{lem:tree}(d), which proves \eqref{eq:phi}.

The edge $e$ is a term of the cut of $S$ in the colour $z$: it joins $x \in S$ to $y \notin S$, and if
$y = O$ then $\xi_x$ is real because $\Psi(e)$ is real. Its value is $\Phi^z(e) = \Ps + h$, the largest
possible by \eqref{eq:phi}. By the cut condition a term $e' \ne e$ has the same value, and by
\eqref{eq:phi} it lies in $\Es$ and has the node $A$. The statement for $\nd{e}{\zeta}$ follows from
the same argument with the two colours exchanged.

For the simple cycle, note that $\Es$ is not empty and that every vertex of $\Gamma(\Es)$ lies on at
least two edges. Take a path of $\Gamma(\Es)$ with distinct vertices and the largest possible number of
edges. Its last vertex lies on an edge that is not the last edge of the path, and the other end of that
edge is a vertex of the path, because the path cannot be made longer. The part of the path from that
vertex to the last one, with the additional edge, is a simple cycle. It has an even number of edges
because $\Gamma(\Es)$ is bipartite, and its edges can be named $f_0, g_0, \dots, f_{r-1}, g_{r-1}$ as in
Definition~\ref{def:alt}.

\step{Step 2: the relation of (c).} Let $f_0, g_0, \dots, f_{r-1}, g_{r-1}$ be a closed alternating
sequence in $\Es$. The equations $\omega_{f_i} = \omega_{g_i}$ and $\zeta_{g_i} = \zeta_{f_{i+1}}$ give
$\sum_e \varepsilon(e) (\zeta_e + \omega_e) = 0$, the sum over the $2r$ edges. On $\Es$,
\[
  \zeta_e + \omega_e = 2\Ps + 2\, \one{e} \cdot u \ \text{ for a body pair } e,
  \qquad
  \zeta_e + \omega_e = -\Ps - u_t \ \text{ for } e = \{t, O\} .
\]
Let $\sigma$ be the sum of $\varepsilon(e)$ over the body pairs of the sequence and $\sigma_O$ the sum
over its $O$-edges. The sum above is $(2\sigma - \sigma_O) \Ps + q \cdot u$. The sequence has $r$ edges
of each sign, which gives $\sigma + \sigma_O = 0$, and therefore $2\sigma - \sigma_O = 3\sigma = N$.
This proves $N \Ps + q \cdot u = 0$. We note for Step~4 that $\sum_j q_j = 4\sigma - \sigma_O = 5N/3$.

\step{Step 3: proof of (b).} Suppose that all edges of a closed alternating sequence in $\Es$ are body
pairs. The restrictions of $\zeta$ and $\omega$ to the bodies are ultrametrics with real values. A body
pair attains $\Ps$, and therefore $\Po = \Ps$ and the pairs of the sequence lie in $\Eo$. The sequence
is then a closed alternating sequence in $\Eo$, which Lemma~\ref{lem:C} excludes because
$\tau \ge 2n^2 \ge n^2$.

\step{Step 4: $N \ne 0$.} By Step~3 the sequence contains an $O$-edge $\{t, O\}$, and it contains it
once, because its edges are distinct. In $q_t$ every body pair through $t$ contributes $\pm 2$ and the
edge $\{t, O\}$ contributes $\pm 1$. Hence $q_t$ is odd and $q \ne 0$. Suppose that $N = 0$. Then
$\sum_j q_j = 0$ and $q \cdot u = 0$ by Step~2. A body lies in at most $n - 1$ body pairs and one
$O$-edge, which gives $|q_j| \le 2n - 1$. The sum of Lemma~\ref{lem:B} for $q$ has at most $n - 2$ terms
and is at most $(n - 2)(2n - 1) < 2n^2 \le \tau$. Lemma~\ref{lem:B} gives $q \cdot u \ne 0$, a
contradiction.

\step{Step 5: proof of (d).} By (a) and (c) there is a closed alternating sequence with
$\Ps = -q \cdot u / N$. Here $N$ is a nonzero multiple of $3$ and $|q_j| \le 2n - 1$, which gives the
bound. The pairs $(N, q)$ of all closed alternating sequences of edges on $n + 1$ vertices form a finite
set that depends on $n$ only.
\end{proof}

\begin{corollary}\label{cor:bound}
Let $n \ge 2$, $\tau \ge 2n^2$ and $u \in W_n(\tau)$. For every weight datum $(u, \zeta, \omega)$ and
every body pair $\{t, l\}$,
\[
  \zeta_{tl} + \omega_{tl} \;\le\; \tfrac23 (2n - 1) \sum_j |u_j| \;+\; 2u_t + 2u_l .
\]
\end{corollary}

\begin{proof}
By Definition~\ref{def:psi}, $\zeta_{tl} + \omega_{tl} = 2\Psi(t, l) + 2u_t + 2u_l$, and
$\Psi(t, l) \le \Ps$. The bound on $\Ps$ is Theorem~\ref{thm:D4}(d).
\end{proof}

\begin{example}\label{ex:n2cycle}
In Example~\ref{ex:n2} the two edges of $\Es$ form the simple cycle $f_0 = \{0, 1\}$, $g_0 = \{1, O\}$,
with $N = -3$ and $q = (-2, -3)$. The relation is $-3\Ps - 2u_0 - 3u_1 = 0$.
\end{example}

\section{Finiteness over the field of Hahn series}\label{sec:K}

In this section a \emph{solution} is a solution of $\cN_n$ in $K$ at fixed masses $m \in (K^{*})^n$, and
the solutions form a Zariski-closed subset of an affine space over the algebraically closed field $K$.
The number $\zeta_{tl} + \omega_{tl}$ is the valuation of $X_{tl} Y_{tl}$, which for a central
configuration is the square of a mutual distance. A large valuation is a close approach of two bodies.

\begin{proposition}\label{prop:D2}
Let $m \in (K^{*})^n$. Suppose that there is a real number $M$ such that
$v(X_{tl} Y_{tl}) \le M$ for every solution of $\cN_n$ in $K$ with the masses $m$ and every body pair
$\{t, l\}$. Then no irreducible component of the set of these solutions has dimension at least $2$.
\end{proposition}

\begin{proof}
Let $V$ be an irreducible component of dimension at least $2$. We show first that every function
$X_{tl} Y_{tl}$ and every $P_{tl}$ is constant on $V$, and then that $V$ contains finitely many orbits
of the rotation, which contradicts its dimension.

\step{Step 1: every $X_{tl} Y_{tl}$ is constant on $V$.} Fix a body pair and let $\phi = X_{tl} Y_{tl}$,
a regular function on $V$. The image $\phi(V)$ is a constructible subset of the affine line, by
Chevalley's theorem, and its closure is irreducible. Hence $\phi$ is constant, or $\phi(V)$ contains
all elements of $K$ except finitely many. In the second case $\phi(V)$ contains an element of valuation
$M + 1$, by the third property of $v$, and this contradicts the hypothesis. Hence $\phi$ is a constant
$\delta_{tl}$, and $\delta_{tl} \ne 0$ by Lemma~\ref{lem:facts}(a).

\step{Step 2: every $P_{tl}$ is constant on $V$.} By \eqref{eq:N1}, $P_{tl}^2 = \delta_{tl}^{-1}$ on
$V$. A regular function on an irreducible variety that takes two values at most is constant.

\step{Step 3: the slice.} The rotation maps $V$ to itself, because $K^{*}$ is connected. Fix two bodies
$a \ne b$ and let $V_1$ be the set of the points of $V$ with $X_{ab} = 1$. The function $X_{ab}$ has no zero on $V$, and the rotation by $s$ multiplies it by $s$. The morphism
$K^{*} \times V_1 \to V$ given by the rotation therefore has the regular inverse that sends a point $p$
of $V$ to the pair of $X_{ab}(p)$ and the image of $p$ under the rotation by $X_{ab}(p)^{-1}$. It is an
isomorphism, and $\dim V_1 = \dim V - 1 \ge 1$.

\step{Step 4: $V_1$ is finite.} On $V_1$ we have $X_{ab} = 1$ and $Y_{ab} = \delta_{ab}$. For $n = 2$
this fixes all differences. When there is a third body $c$, we have $X_{cb} = X_{ca} + 1$ and
$Y_{cb} = Y_{ca} + \delta_{ab}$, and therefore
\[
  X_{ca} Y_{ca} = \delta_{ca} \qquad\text{and}\qquad
  (X_{ca} + 1)(Y_{ca} + \delta_{ab}) = \delta_{cb} .
\]
The second equation minus the first gives
$Y_{ca} = \delta_{cb} - \delta_{ca} - \delta_{ab} - \delta_{ab} X_{ca}$, and the first then reads
\[
  -\delta_{ab}\, X_{ca}^2 + (\delta_{cb} - \delta_{ca} - \delta_{ab})\, X_{ca} - \delta_{ca} = 0 .
\]
Its leading coefficient is not $0$; a quadratic equation for the configurations with given mutual
distances is used in the same way in~\cite[\S 2.1]{tomar6}. Hence $X_{ca}$ and $Y_{ca}$ take at most
two values each on $V_1$,
and every difference $X_{tl}$, $Y_{tl}$ takes finitely many values on $V_1$. By Step~2 and the equations
\eqref{eq:N2} and \eqref{eq:N3}, the coordinates $X_t$ and $Y_t$ take finitely many values on $V_1$.
The set $V_1$ is therefore finite, which contradicts $\dim V_1 \ge 1$.
\end{proof}

The proposition uses no balancing condition and no theorem on the dimension of a tropical variety. A
component on which bodies only separate, or on which one position tends to $0$, is excluded with the
others, because on a component of dimension at least $2$ some $X_{tl} Y_{tl}$ takes values of every valuation.

\begin{proof}[Proof of Theorem~\ref{thm:K}]
Let $u_t = v(m_t)$, so that $u \in W_n(2n^2)$. By Lemma~\ref{lem:A} the weights of a solution form a
weight datum $(u, \zeta, \omega)$. By Corollary~\ref{cor:bound}, with $\tau = 2n^2$, every solution and
every body pair satisfy
$v(X_{tl} Y_{tl}) = \zeta_{tl} + \omega_{tl} \le \tfrac23 (2n - 1) \sum_j |u_j| + 4 \max_j |u_j|$. By
Proposition~\ref{prop:D2} no irreducible component of the set of solutions has dimension at least $2$,
and by Lemma~\ref{lem:orbits} the solutions form finitely many orbits of the rotation.
\end{proof}

\section{Generic masses}\label{sec:generic}

Let $R = \QQ[m_0, \dots, m_{n-1}]$ and $\cM = \Spec R$, and let $\pi \colon \cX \to \cM$ be the affine
scheme of finite presentation over $\cM$ defined by the equations \eqref{eq:N1} to \eqref{eq:N3} in the
unknowns $X_t$, $Y_t$, $P_{tl}$. For a field $k$ of characteristic $0$, a vector $m \in k^n$ is a
morphism $\Spec k \to \cM$, and its image is the \emph{point of $\cM$ below $m$}. The scheme of the
solutions at the masses $m$ is the fibre of $\cX$ at that point, extended to $k$. For a point $x$ of
$\cX$, the \emph{fibre dimension at $x$}, written $\dim_x \cX_{\pi(x)}$, is the largest dimension of an
irreducible component through $x$ of the fibre of $\pi$ that contains $x$. If $k$ is algebraically
closed, the closed points of the scheme of the solutions are the solutions in $k$. They are dense in
every irreducible component of the scheme, and the irreducible components of the scheme and of the
Zariski-closed set of the solutions in $k$ correspond, with the same dimensions; nilpotent elements of
the coordinate ring do not affect them.

\begin{lemma}\label{lem:D0}
There is a constructible subset $D$ of $\cM$ with the following two properties.
\begin{enumerate}
\item[(a)] For every algebraically closed field $k$ of characteristic $0$ and every $m \in k^n$, the
system $\cN_n$ with the masses $m$ has infinitely many solutions in $k$ modulo the rotation if and only
if the point of $\cM$ below $m$ lies in $D$.
\item[(b)] If $D$ is dense in $\cM$, then for every $u^{*} \in \RR^n$ whose coordinates are linearly
independent over $\QQ$, the point of $\cM$ below the vector
$(T^{u^{*}_0}, \dots, T^{u^{*}_{n-1}}) \in K^n$ lies in $D$.
\end{enumerate}
\end{lemma}

\begin{proof}
The set $\cX_{\ge 2} = \{ x \in \cX : \dim_x \cX_{\pi(x)} \ge 2 \}$ is closed in
$\cX$~\cite[Tag~02FZ]{stacks}; this is Chevalley's semicontinuity theorem \cite[13.1.3]{ega4iii}. Let
$D = \pi(\cX_{\ge 2})$. The morphism $\pi$ is affine and of finite presentation, and $\cX_{\ge 2}$ is a
closed subset of a Noetherian scheme; therefore $D$ is constructible~\cite[Tag~054K]{stacks}.

(a) Let $s$ be the point of $\cM$ below $m$, with residue field $\kappa(s) \subseteq k$. The scheme of
the solutions over $k$ is $\cX_s \otimes_{\kappa(s)} k$. The fibre dimension at a point does not
change under an extension of the base field, and every point of $\cX_s$ is the image of a point of
$\cX_s \otimes_{\kappa(s)} k$~\cite[Tag~02FY, parts (1) and (3)]{stacks}. Hence the scheme of the
solutions over $k$ has a point at which its dimension is at least $2$ if and only if $s \in D$. Such a
point exists if and only if an irreducible component of the scheme has dimension at least $2$. By the
correspondence stated before the lemma, this holds if and only if an irreducible component of the set
of the solutions in $k$ has dimension at least $2$, and by Lemma~\ref{lem:orbits} if and only if there
are infinitely many orbits.

(b) The scheme $\cM$ is irreducible, and a dense constructible subset of an irreducible space contains
its generic point~\cite[Tag~005K]{stacks}. The residue field of the generic point is
$F = \QQ(m_0, \dots, m_{n-1})$. The ring homomorphism $R \to K$ with $m_t \mapsto T^{u^{*}_t}$ sends a
polynomial $\sum_{\alpha} h_\alpha m^{\alpha}$ to $\sum_{\alpha} h_\alpha T^{u^{*} \cdot \alpha}$. The
exponents $u^{*} \cdot \alpha$ are distinct for distinct $\alpha$, because the coordinates of $u^{*}$ are
linearly independent over $\QQ$. The homomorphism is therefore injective, and the point of $\cM$ below
$(T^{u^{*}_t})$ is the generic point.
\end{proof}

\begin{proof}[Proof of Theorem~\ref{thm:Nn}]
The set $W_n(2n^2)$ is open and not empty, and hence it contains a vector $u^{*}$ whose coordinates are
linearly independent over $\QQ$. By Theorem~\ref{thm:K} the system $\cN_n$ with the masses
$m_t = T^{u^{*}_t}$ has finitely many solutions in $K$ modulo the rotation. By Lemma~\ref{lem:D0}(a) the
point of $\cM$ below these masses does not lie in $D$, and by Lemma~\ref{lem:D0}(b) the set $D$ is not
dense in $\cM$. There is therefore a nonzero polynomial $H \in R$ that vanishes on $D$. Let
$\Sigma_n = \{ m \in \CC^n : H(m) = 0 \}$. If $\cN_n$ has infinitely many solutions in $\CC$ modulo the
rotation at $m \in \CC^n$, then the point of $\cM$ below $m$ lies in $D$ by Lemma~\ref{lem:D0}(a), and
$H(m) = 0$.
\end{proof}

\begin{proof}[Proof of Theorem~\ref{thm:main}]
Let $\Sigma_n$ be the set of Theorem~\ref{thm:Nn} and let $m \in \RR_{>0}^n$ lie outside it. The system
$\cN_n$ has finitely many solutions in $\CC$ modulo the rotation, and Lemma~\ref{lem:cc-sol}(c) gives the
statement.
\end{proof}

\begin{remark}\label{rem:roberts}
Theorem~\ref{thm:K} is stronger than what the proof of Theorem~\ref{thm:Nn} uses: it holds for all masses
in $K$ with valuations in $W_n(2n^2)$, and, after a permutation of the bodies, for valuations in the
image of $W_n(2n^2)$ under a permutation of the coordinates. This restricts the closure $\Sigma$ of the
set $D$ of Lemma~\ref{lem:D0} in $\cM$, as follows. Because $D$ is constructible, it contains an open
subset $U$ of $\Sigma$ that is dense in every irreducible component of $\Sigma$. Extend $\Sigma$ and
$U$ to $K$ and consider the part where no mass is $0$, a closed subset of a torus over $K$. Each of its
irreducible components lies over a component of $\Sigma$ and meets the extension of $U$, because the
extension of scalars is flat and therefore maps generic points of components to generic points of
components. Suppose that a
point over $K$ of such a component has its vector of valuations in one of the open cones. The points
over $K$ of the component with that vector of valuations are then Zariski dense in
it~\cite[Corollary~4.2, and Remark~2 of the erratum]{payne2009}, and one of them lies in the extension
of $U$; at that point Lemma~\ref{lem:D0}(a) contradicts Theorem~\ref{thm:K}. Hence no point of $\Sigma$
over $K$ with nonzero masses has its vector of valuations in these open cones. A mass vector over $K$, or
over a valued field with an embedding in $K$ that preserves the valuation, at which $\cN_n$ has infinitely
many solutions modulo the rotation and whose valuations lie in $W_n(2n^2)$ would contradict the theorem.
The examples of continua that the search of 17 November 2026 found (Section~\ref{sec:intro}) have
equal valuations: Roberts' continuum for five
bodies~\cite{roberts1999} has the masses $(1, 1, 1, 1, -1/4)$ up to a common factor. At the positive
masses $(4, 4, 4, 4, 1)$ the system $\cN_5$ has a continuum of solutions in which $P_{t4}$ is the
negative root for the four pairs that contain the light body
\cite[p.~543 and Remark~8]{albouykaloshin2012}; these solutions are not central configurations, and
Tomar~\cite{tomarspecial} finds 36 curves of them. If $n - 1$ masses are $0$ and $n \ge 3$, the
massless bodies move freely on a circle about the heavy body; such masses are not in $(K^{*})^n$.
\end{remark}

\section{Where the hypotheses are used}\label{sec:hyp}

Table~\ref{tab:hyp} lists, for each feature of the system and each hypothesis of Theorem~\ref{thm:K}, the
step that uses it. The positivity of the masses is used in Lemma~\ref{lem:norm} only, and no proof uses an identity for the
angular momentum of a cluster.

\begin{table}[ht]
\caption{The features of the system $\cN_n$ and the hypothesis on the mass weights, with the steps that
use them.}\label{tab:hyp}
\small
\begin{tabular}{@{}>{\raggedright}p{0.36\textwidth}>{\raggedright\arraybackslash}p{0.58\textwidth}@{}}
\toprule
feature or hypothesis & where it is used \\
\midrule
the exponent $3$ of $P_{tl}$ in \eqref{eq:N2}, \eqref{eq:N3} & Lemma~\ref{lem:A} (the value $\Phi^z = \Psi + \omega$); Theorem~\ref{thm:D4}, Step~2 (the factor $3$ of $N$) and Step~4 (parity: a body pair contributes an even number to $q_t$ and an $O$-edge an odd one) \\
the left sides $X_t$, $Y_t$ of \eqref{eq:N2}, \eqref{eq:N3} & the $O$-edges. With left sides $\theta X_t$, $\theta Y_t$ and a further unknown $\theta$, Step~2 has the additional term $(N/3)\, v(\theta)$, $\Ps$ is not bounded, and the system has a continuum by scaling \\
every pair of bodies interacts & Corollary~\ref{cor:bound}: every body pair is an edge, and $\Psi \le \Ps$ bounds the valuation of $X_{tl} Y_{tl}$ \\
$u \in W_n(\tau)$ with $\tau \ge n^2$ & Lemma~\ref{lem:C}, Steps~1 and~4; Example~\ref{ex:0011} shows that it is needed there \\
$u \in W_n(\tau)$ with $\tau \ge 2n^2$ & Theorem~\ref{thm:D4}, Step~4 \\
\bottomrule
\end{tabular}
\end{table}

\begin{remark}\label{rem:u0123}
We do not know whether Theorem~\ref{thm:D4}(c) holds for every vector $u$ with distinct coordinates. A
floating-point search with linear programs (Section~\ref{sec:comp}) at $u = (0, 1, 2, 3)$ finds triples
that satisfy the conditions of a weight datum within the tolerance of the search and have the closed
alternating sequence
$\{0, 2\}, \{2, O\}, \{0, O\}, \{0, 3\}$ with $N = 0$ and $q = (1, 0, -3, 2)$, and with values of $\Ps$
up to the largest value that the search admits. Lemma~\ref{lem:B} excludes this vector $q$ in $W_4(\tau)$ for
$\tau \ge 1/2$. This is evidence that the hypothesis on $u$ cannot be dropped from
Theorem~\ref{thm:D4}; it says nothing about Theorem~\ref{thm:K}, because a weight datum need not come
from a solution.
\end{remark}

\begin{remark}\label{rem:exponent}
For a force with the exponent $\varrho$ in place of $3$ in \eqref{eq:N2} and \eqref{eq:N3}, the term
of a body pair in a cut identity is $m_t m_l P_{tl}^{\varrho} X_{tl}$, and the value that takes the
place of $\Psi(t, l)$ is $\tfrac12 (\varrho - 2)(\zeta_{tl} + \omega_{tl}) - u_t - u_l$. For
$\varrho < 2$ the coefficient of $\zeta_{tl} + \omega_{tl}$ is negative. The triangle conditions of
Step~4 of Lemma~\ref{lem:C} then bound $\Psi(x, x')$ from above, as the maximality of $\Po$ does, and
the step gives no contradiction. Continua at an open set of masses exist for such forces: for the
potentials $\sum m_t m_l\, r_{tl}^{2i}$ with an integer $i \ge 2$, which have $\varrho = 2 - 2i$,
Nurowski~\cite{nurowski2026} gives one-parameter families of central configurations of $2(i + 1)$
bodies that persist for all masses near the equal masses. In these
families two bodies collide for some values of the parameter, and the force, a polynomial, vanishes
there.
\end{remark}

\section{The formal proof}\label{sec:lean}

The files named in this section and in Sections~\ref{sec:comp} and~\ref{sec:llm} are in the repository
of this paper, \url{\repourl}, and their paths are relative to its root.
All theorems, propositions, lemmas and corollaries of Sections~\ref{sec:system} to~\ref{sec:generic} are
proved in the text. Table~\ref{tab:lean} lists those for which, in addition, a formal statement is
proved in Lean~4~\cite{lean4} with the library Mathlib~\cite{mathlib}, and those for which there is
none.

\begin{table}[ht]
\caption{The statements of this paper and their formal counterparts. The formal statements are in the
files \repo{Challenge.lean} (280 lines), \repo{ChallengeA.lean} (57 lines) and \repo{ChallengeB.lean}
(103 lines) of the directory \repo{lean/direct-route} of the repository of this paper.}\label{tab:lean}
\small
\begin{tabular}{@{}>{\raggedright}p{0.27\textwidth}>{\raggedright}p{0.27\textwidth}>{\raggedright\arraybackslash}p{0.39\textwidth}@{}}
\toprule
statement & formal name & what the formal statement covers \\
\midrule
Lemmas~\ref{lem:norm} to~\ref{lem:orbits} & none & not formalised \\
Lemma~\ref{lem:tree} & none & not formalised as a statement \\
Lemma~\ref{lem:A} & \repo{lemmaA} & the last sentence: the weights of a solution with nonzero masses form a weight datum, for every field with a valuation in $\RR \cup \{+\infty\}$. The identities \eqref{eq:cut} and the valuations of their terms are not formal statements \\
Lemma~\ref{lem:B} & \repo{lemmaB}, \repo{lemmaB_sign} & the whole statement \\
Definition~\ref{def:alt} & \repo{IsAltSeq}, \repo{IsCycle} & definitions \\
Lemma~\ref{lem:C} & \repo{lemmaC_alt} & the whole statement \\
Corollary~\ref{cor:C} & \repo{lemmaC}, \repo{cycle_alt_bodies} & the whole statement, and that a simple cycle is a closed alternating sequence \\
Theorem~\ref{thm:D4}(a) & \repo{degree_two}, \repo{exists_cycle} & the whole statement, for every weight datum, with no hypothesis on $u$ \\
Theorem~\ref{thm:D4}(b), (c) & \repo{alt_relation}, \repo{cycle_relation}, \repo{cycle_alt} & the whole statements; \repo{cycle_relation} is the case of the simple cycles \\
Theorem~\ref{thm:D4}(d) & \repo{psi_star_bound} & the inequality; that $\Ps$ lies in a finite set is not a formal statement \\
Corollary~\ref{cor:bound} & \repo{collision_bound}, \repo{solution_bound} & the bound for a weight datum, and the same bound on $v(X_{tl} Y_{tl})$ at a solution over a valued field \\
Examples~\ref{ex:n2} and~\ref{ex:n2cycle} & \repo{ex2_spec}, \repo{ex2_alt} & the weight datum at $u = (0, 2)$, its set $\Es$, its cycle, $N$ and $q$; the equation for $P_{01}$ and the formulas for a general $u$ are not formal statements \\
Example~\ref{ex:0011} & \repo{ex4_cycle}, \repo{ex4_alt} & the two ultrametric conditions and the displayed cycle at the stated weights \\
$W_n(\tau) \ne \emptyset$ & \repo{chamber_nonempty} & for all $n$ and all $\tau \ge 0$ \\
Proposition~\ref{prop:D2}, Lemma~\ref{lem:D0} & none & not formalised \\
Theorems~\ref{thm:main}, \ref{thm:Nn}, \ref{thm:K} & none & not formalised \\
\bottomrule
\end{tabular}
\end{table}

\subsection*{Versions and the check} The formal proof uses Lean~4.36.0 and Mathlib at the tag
\repo{v4.36.0}. It is checked with \repo{lake comparator}, the tool distributed with Lean that compares
a file of statements, in which every proof is the placeholder \repo{sorry}, with a file of proofs: it
accepts if every listed theorem is proved with the statement of the first file and the proof uses no
axiom outside a permitted list. The permitted axioms are \repo{propext}, \repo{Quot.sound} and
\repo{Classical.choice}. The three statement files use \repo{sorry} as the proof of each stated theorem. The
accepted proofs of the listed theorems do not depend on the axiom \repo{sorryAx}, directly or through
another declaration, because the list of permitted axioms excludes it. With the option \repo{--paranoid} the proof is accepted by the kernel of Lean
and by the five further checkers that the tool runs (\repo{leanchecker-paranoid}, \repo{lean4lean},
\repo{nanoda}, \repo{con-leche}, \repo{con-ron}), on an ARM machine under macOS and on an x86-64 machine under
Linux. The tool can run the proof files inside a sandbox; on
both machines it ran without the sandbox.

\subsection*{The statement files and their reviews} Each of the three statement files was compared
with the text by a separate language-model session that had the file and the text and no formal proof
(Section~\ref{sec:llm}); the three reports are in the repository. The file \repo{Challenge.lean} holds
the definitions and the statements of Sections~\ref{sec:chamber} to~\ref{sec:D4} for simple cycles. The
file \repo{ChallengeA.lean} states Lemma~\ref{lem:A} and the form of Corollary~\ref{cor:bound} at a
solution. The file \repo{ChallengeB.lean} states Lemma~\ref{lem:B} with the sign, Lemma~\ref{lem:C}, and
Theorem~\ref{thm:D4}(b), (c) and~(d) in the form printed here. No review asked for a change to the
statement of a theorem; the first proposed two further examples, which were added. The three files are
fixed, and Table~\ref{tab:sha} gives their SHA-256 digests. The file
\repo{lean/direct-route/TRUSTED.sha256} holds these digests and those of the lists of theorems and of
the files that select the versions, and the script \repo{scripts/allN/lean_gate.sh} checks all of them
before it runs the comparison. Two controls were run: with a \repo{sorry} placed in one
proof the check fails, and with a false hypothesis added to a statement the check fails because the
statement file differs from the fixed one.

\begin{table}[ht]
\caption{SHA-256 digests of the three files of formal statements.}\label{tab:sha}
\footnotesize
\begin{tabular}{@{}ll@{}}
\toprule
file & SHA-256 \\
\midrule
\repo{Challenge.lean} & \repo{98bc0a4f29f32f26484136717c7ac15b05efd387e3344a048fbbe0b0cf48feec} \\
\repo{ChallengeA.lean} & \repo{eb0bceb664c95a884e4f8a5453062c77279ff8fe9db45cd1d9beb41f1c697b22} \\
\repo{ChallengeB.lean} & \repo{1879a2be79d0963d85317940e82ff2d98ca3a3421ea56a2b51000d615a971038} \\
\bottomrule
\end{tabular}
\end{table}

\subsection*{Differences between the formal statements and the text} The formal statements differ from
those of the text in the following ways.
\begin{enumerate}
\item The bodies are the natural numbers below $n$ and the vertex $O$ is the number $n$. Vectors are
functions on the natural numbers. For $n \ge 2$ only the values at the indices below $n$ are used.
\item The values $\xi_t$ and $\eta_t$ lie in the type of the real numbers with a top element, and the
top element stands for $+\infty$. An $O$-edge $\{t, O\}$ is admitted as an edge only if $\xi_t$ and
$\eta_t$ are both real; in the text an $O$-edge with an infinite value has $\Psi = -\infty$ and never
lies in $\Es$. A term $\{t, O\}$ of a cut needs only the value of its own colour to be real, as in
Definition~\ref{def:psi}.
\item The number $\Ps$ is a parameter with the hypothesis that it bounds $\Psi$ on the edges and that
an edge attains it.
\item A simple cycle or a closed alternating sequence is given by two functions $i \mapsto f_i$ and
$i \mapsto g_i$. That the $2r$ edges are distinct is written as three conditions: the $f_i$ are
distinct, the $g_i$ are distinct, and no $f_i$ equals a $g_j$.
\item Lemma~\ref{lem:B} is stated with sums over the indices up to a given index $p$ with $q_p \ne 0$,
in place of the largest index with a nonzero coefficient. The statement on the sign says that the
product of $q_p$ and $q \cdot u$ is positive.
\item Theorem~\ref{thm:D4}(b) and~(c) form one formal statement.
\item The field of Lemma~\ref{lem:A} is any field with a valuation in the real numbers with a top
element; $K$ is one such field.
\item The formal counterparts of Theorem~\ref{thm:D4}(a) have no hypothesis on $u$; the proof of
Step~1 uses none.
\end{enumerate}
The reviews list two further differences of encoding. The diagonal entries of $\zeta$ and $\omega$ are
stored, and no statement uses them. For $n < 2$, which the text excludes, the formal definition of
$W_n(\tau)$ consists of the condition $u_0 < u_1$ on a padded vector, and \repo{chamber_nonempty} holds in
that form.

\subsection*{The hypotheses can hold} A statement whose hypotheses nothing satisfies is true and says
nothing. The file \repo{Challenge.lean} therefore contains, as theorems that are proved with the
others, a vector in $W_n(\tau)$ for all $n$ and all $\tau \ge 0$, the weight datum of Example~\ref{ex:n2} with its
cycle, a weight datum for three bodies over the vector $(0, 1, 20)$ of $W_3(18)$, and two weight data in
which a position has the valuation $+\infty$. The file \repo{ChallengeB.lean} contains the closed
alternating sequences of Examples~\ref{ex:n2cycle} and~\ref{ex:0011}. The second review gives a solution of $\cN_2$ over $K$
with mass weights $(3, 4)$, which satisfies the hypotheses of Lemma~\ref{lem:A} and of the form
of Corollary~\ref{cor:bound} at a solution; this solution is not formalised.

\section{Computations}\label{sec:comp}

The searches of Table~\ref{tab:comp} are in floating point, with no certificates, and concern
Theorem~\ref{thm:D4}. Two exact enumerations at the end of the section concern Lemma~\ref{lem:C}. All
of them are evidence, and no proof uses them. The program \repo{scripts/allN/nk_fibre.py} of the
repository takes a vector $u$ and, for each ordered pair of binary trees on the
$n + 1$ vertices, searches with linear programs (HiGHS~\cite{highs}) for heights $\zeta$, $\omega$ on
these trees that satisfy every cut condition. A tie between two terms is accepted within the absolute
tolerance $10^{-3}$. A second program, \repo{nk_cut_check.py}, tests the ultrametric conditions and the
cut conditions of each recorded point without the trees and the linear programs, and a third,
\repo{nk_cycle_check.py}, builds $\Es$ and $\Gamma(\Es)$ at each point and evaluates
$N\Ps + q \cdot u$ on a basis of the cycles.

\begin{table}[ht]
\caption{Floating-point searches for weight data. ``Pairs of trees'' is the number of ordered pairs
of binary trees searched and its share of all pairs. A threshold $R$ restricts the search to points
with $\Ps \ge R$. The last column is the bound of Theorem~\ref{thm:D4}(d). The vectors marked
with a dagger lie in no $W_n(\tau)$ with $\tau \ge (n - 2)(2n - 1) - 1$.}\label{tab:comp}
\footnotesize
\setlength{\tabcolsep}{4pt}
\begin{tabular}{@{}l>{\raggedright}p{0.27\textwidth}>{\raggedright}p{0.13\textwidth}lrrr@{}}
\toprule
$n$ & $u$ & pairs of trees & $R$ & points & largest $\Ps$ & bound \\
\midrule
5 & $(0, 1, 12, 150, 1900)^{\dagger}$ & 893,025 (all) & 5 & 203 & 2,424.3 & 6,189 \\
5 & $(0, 1, 2.37, 4.91, 8.53)^{\dagger}$ & 893,025 (all) & 3 & 30 & 5.52 & 50.4 \\
5 & $(0, 1, 52, 2700, 140000)$ & 893,025 (all) & 5 & 204 & 184,831 & 428,259 \\
6 & $(0, 1, 45, 2000, 89000, 3920000)$ & 100,000 (0.09\%) & none & 1,589 & 2,613,303 & 14,707,169 \\
6 & $(0, 1, 45, 2000, 89000, 3920000)$ & 500,000 (0.46\%) & 5 & 189 & 7,749,636 & 14,707,169 \\
6 & $(0, 1, 12, 150, 1900, 24000)^{\dagger}$ & 300,000 (0.28\%) & 5 & 108 & 31,837 & 95,564 \\
\bottomrule
\end{tabular}
\end{table}

Table~\ref{tab:comp} lists the searches for five and six bodies. At $n = 6$ the pairs of trees are a
random sample of the 108,056,025 ordered pairs, drawn with the seed $29$. All $2{,}323$ points of the
table pass the second program. At each of them every vertex of $\Gamma(\Es)$ lies on two edges, every
cycle of the basis has $N \ne 0$, and $|N\Ps + q \cdot u| \le 2 \cdot 10^{-5}$; no point lies above the
bound of the last column. The vector of the third row lies in $W_5(50)$, as Theorem~\ref{thm:D4} requires for $n = 5$.
The vector of the fourth and fifth rows lies in $W_6(43)$ and not in $W_6(72)$. The proofs of
Lemma~\ref{lem:C} and Theorem~\ref{thm:D4} use the hypothesis on $\tau$ only through the inequalities
$\tau \ge (n - 2)(n + 1) - 1$ and $\tau \ge (n - 2)(2n - 1) - 1$, which read $\tau \ge 27$ and
$\tau \ge 43$ for $n = 6$. At the three vectors with a dagger the theorem makes no statement, and the
relation and the bound hold at every point found.

The searches find a failure where there is one. With the threshold on $\Ps$ at $n = 3$ and
$u = (0, 1, 30)$, where every weight datum has $\Ps = -1$, the search finds points for the thresholds
$-40$, $-20$ and $-1.5$ and none for $-0.9$, $0$, $10$ and $1000$. The points of Remark~\ref{rem:u0123}, with $N = 0$, were
found at $u = (0, 1, 2, 3)$ by a separate program with the same method
(\repo{data/n7/allN/hier/r29/r29_batch.py}); the same program finds no such point at the vectors
$(0, 1, 30, 900)$, $(0, 1, 4, 16)$ and $(0, 1, 2.37, 4.91)$.

Steps~2 to~4 of the proof of Lemma~\ref{lem:C} are not tested by these searches below six bodies. An
exact enumeration of all pairs of trees
(the program \repo{search_balanced_cycles.py} in \repo{data/n7/allN/hier/r31} of the repository) finds, for $n \le 5$, no closed alternating
sequence with $q = 0$ in which every pair consists of the two heaviest bodies of the intersection of
its two nodes, a property that every pair of $\Eo$ has when the mass weights are distinct. It finds $96$
such sequences for six bodies, each through all six, and we ran no search for $n \ge 7$.

A second exact search tests Lemma~\ref{lem:C} itself at one vector for each $n$ from $2$ to $6$, namely
$u_0 = 0$ and $u_j = (n^2 + 2)^{j-1}$ for $j \ge 1$, which lies in $W_n(n^2)$. The program
\repo{search_alternating.cpp} in \repo{data/n7/allN/hier/r32} of the repository, written by the session
that refereed this paper (Section~\ref{sec:llm}), enumerates for each colour the order types of the
heights of an ultrametric on $n$ bodies, of which there are $9{,}012$ for $n = 6$. For each unordered
pair of order types, $40{,}612{,}578$ for $n = 6$, it decides in integer arithmetic whether real heights
of these types exist for which $\Eo$ contains a closed alternating sequence. It finds no such
heights for $n \le 6$, and at $u = (0, 0, 1, 1)$ it finds $82$ feasible simple cycles in its
graphs of height ranks, counted over the pairs of order types, among them the cycle of
Example~\ref{ex:0011}. We ran the program again and obtained the same output, and a second program,
written independently by another session (\repo{check_rank_cycles.cpp} in \repo{data/n7/allN/hier/r33}),
reproduces all counts.

\section{Problems}\label{sec:problems}

\begin{problem}
Describe a set $\Sigma_n$ for which Theorem~\ref{thm:Nn} holds, or give a bound for the number of
solutions of $\cN_n$ modulo the rotation, for masses outside it, that is sharper than the bound of
Remark~\ref{rem:scope}.
\end{problem}

\begin{problem}
For which exponents $\varrho > 2$ of the force does the argument of this paper give finiteness outside
a proper algebraic subset of the mass space? Lemma~\ref{lem:C} uses only that $\varrho - 2$ is positive,
whereas Step~4 of Theorem~\ref{thm:D4} uses a parity that holds for $\varrho = 3$ and fails when
$\varrho$ is an even integer.
\end{problem}

\begin{problem}
Formalise Proposition~\ref{prop:D2} and Lemma~\ref{lem:D0}, and with them Theorem~\ref{thm:Nn}.
\end{problem}

\begin{problem}
Does Theorem~\ref{thm:K} hold for all masses in $K$ whose valuations are distinct, or whose valuations
are linearly independent over $\QQ$?
\end{problem}

\section{Use of language models}\label{sec:llm}

This section records how the results of this paper were produced, which a reader needs in order to
judge them.

\subsection*{Division of labour} The author set the goal, the finiteness theorem outside a proper
algebraic subset for every $n$, after a program on seven bodies that followed the six-body method of
Tomar~\cite{tomar6}. The author also fixed the budget for the computations, decided which checks to
run, and started the separate sessions described below. The mathematics was produced by a
language-model agent (Claude Fable~5.2, Anthropic) with access to a terminal, to the repository and to
rented machines. In working sessions that the author directed, the agent found and wrote the proofs of
Sections~\ref{sec:val} to~\ref{sec:generic}, wrote and ran the programs of Section~\ref{sec:comp},
wrote the two files of formal statements and the formal proofs of Section~\ref{sec:lean}, and wrote the
text of this paper from the documents of the repository. At the time of writing, no person who works
on central configurations or on algebraic geometry has read the proofs.

\subsection*{Separate sessions} Eight further sessions, each started by the author with no access to
the working sessions, checked the work. The first used Qwen4 Max (Alibaba); the other seven used
GPT-6.2 Sol (OpenAI) at its highest setting of reasoning effort. Three sessions read earlier write-ups
of the proof and wrote their own programs to test it. The texts that these three sessions and the
two readers of this paper judged are kept as they were read, in the directory \repo{docs/as-read} of
the repository.
Each of the three reported that it found no error and asked for corrections of wording and of detail,
which were made. The third searched the literature and found that the survey of the working agent had
missed the two six-body preprints~\cite{tomar6,chentong6} and the two claims for every
$n$~\cite{tibboel2013,okolo2026}; Table~\ref{tab:known} follows the corrected survey. Three sessions
compared the three files of formal statements with the text, without access to the formal proofs
(Section~\ref{sec:lean}); none asked for a change to a statement.

Two sessions, run in OpenAI Codex, read drafts of this paper with the instruction to judge the printed
proofs alone. The first read the draft \repo{docs/as-read/R32-paper.tex}. It found that the proof of
Lemma~\ref{lem:tree}(d) cited part~(a) outside its hypotheses; the proof printed here is the repaired
one. It reported no other error in the argument, asked for corrections to the introduction, to
Table~\ref{tab:known}, to Remark~\ref{rem:roberts} and to Section~\ref{sec:lean}, which were made, and
wrote the second exact search of Section~\ref{sec:comp}. The second read the changes made after that
report, in the draft \repo{docs/as-read/R33-paper.tex}. It found the repair and the other changes correct, asked for
further corrections of wording and of scope, which were made, and reproduced the exact search with a
program of its own. These eight reads are not peer review. Their reports are in the directory \repo{docs} of the
repository: the files \repo{n7-referee-R29.md} to \repo{n7-referee-R33.md}, and the three files
\repo{allN-direct-route-lean-review-1.md} to \repo{allN-direct-route-lean-review-3.md}.

\subsection*{What the reads and the formal proof cover in this text} The proof of Lemma~\ref{lem:C}
follows the formal proof and replaces a longer argument that the first three sessions read; the
statement on the sign in Lemma~\ref{lem:B} and the form of Theorem~\ref{thm:D4} for closed alternating
sequences were written with it. Lemma~\ref{lem:D0} is a restatement, with a shorter proof, of the lemma
that those sessions read. The two sessions that read drafts of this paper judged these proofs as
printed in those two drafts. After the second report the text changed in
the places that the report names: wording in the abstract, the introduction, Remark~\ref{rem:scope} and
Remark~\ref{rem:roberts}, and Sections~\ref{sec:lean} to~\ref{sec:llm}; no separate session has read
these last changes. For the statements of Table~\ref{tab:lean} that have a formal counterpart, the
proofs are checked by the kernel of Lean and five further checkers, and not by a language model. What
remains to be trusted there is that the three statement files, which are short, say what the text
says. For Lemmas~\ref{lem:norm} to~\ref{lem:orbits}, Proposition~\ref{prop:D2} and Lemma~\ref{lem:D0},
the text of this paper is the only proof.

\bibliographystyle{amsplain}
\bibliography{references}

\end{document}
